% =====================================================================
%  PREDLOŽAK ZA ZAVRŠNI I DIPLOMSKI RAD  —  LaTeX inačica
%  Sve je u jednoj datoteci. Prevedi s:
%      pdflatex predlozak-zavrsni-diplomski
%      bibtex   predlozak-zavrsni-diplomski
%      pdflatex predlozak-zavrsni-diplomski
%      pdflatex predlozak-zavrsni-diplomski
%  (Na Overleafu je dovoljno pritisnuti "Recompile" 2–3 puta.)
%
%  NE MIJENJAJ preambulu osim onoga što je označeno s  <<< UREDI >>>.
% =====================================================================
\documentclass[12pt,a4paper,oneside]{report}

% ---------------------------------------------------------------------
% <<< UREDI >>>  PODACI O RADU  — mijenjaj samo desnu stranu u vitičastim zagradama
% ---------------------------------------------------------------------
\newcommand{\Sveuciliste}{Sveučilište Jurja Dobrile u Puli}
\newcommand{\Fakultet}{Fakultet informatike u Puli}
\newcommand{\Odjel}{Naziv studija}
\newcommand{\VrstaRada}{DIPLOMSKI RAD}          % ZAVRŠNI RAD / DIPLOMSKI RAD
\newcommand{\NaslovHR}{Naslov rada na hrvatskom jeziku}
\newcommand{\NaslovEN}{Title of the thesis in English}
\newcommand{\NaslovKratki}{Skraćeni naslov rada}  % ide u zaglavlje stranice
\newcommand{\Student}{Ime Prezime}
\newcommand{\JMBAG}{JMBAG: 0000000000}
\newcommand{\Mentor}{doc. dr. sc. Sandi Baressi Šegota}
\newcommand{\Komentor}{}                         % ostavi prazno ako ga nema
\newcommand{\MjestoGodina}{Pula, rujan 2026.}

% ---------------------------------------------------------------------
%  PAKETI
% ---------------------------------------------------------------------
\usepackage[T1]{fontenc}
\usepackage[utf8]{inputenc}
\usepackage{mathptmx}                 % Times (parnjak Times New Romanu iz DOCX-a)
\usepackage[scaled=0.92]{helvet}
\usepackage{courier}
\usepackage{microtype}

% Hrvatski jezik: koristi se ako je paket instaliran (Overleaf, TeX Live full).
% Ako nije, dokument se svejedno prevodi — samo bez hrvatskog prijeloma riječi.
\IfFileExists{croatian.ldf}{\usepackage[croatian]{babel}}{}

\usepackage[a4paper,left=3cm,right=2.5cm,top=2.5cm,bottom=2.5cm]{geometry}
\usepackage[table]{xcolor}
\usepackage{graphicx}
\usepackage{setspace}
\usepackage{titlesec}
\usepackage{titletoc}
\usepackage{fancyhdr}
\usepackage{caption}
\usepackage{booktabs}
\usepackage{array}
\usepackage{tabularx}
\usepackage{enumitem}
\usepackage{amsmath,amssymb}
\usepackage{listings}
\usepackage[numbers,sort&compress]{natbib}
\usepackage{float}
\usepackage[hidelinks]{hyperref}
\usepackage{xurl}

% ---------------------------------------------------------------------
%  BOJE  (Hawaiian blue = službena boja fakulteta)
% ---------------------------------------------------------------------
\definecolor{fakplava}{HTML}{00C3E3}   % akcentna boja: linije, brojevi, zaglavlje tablice
\definecolor{fakplavatamna}{HTML}{00788C} % tamnija inačica za sitan tekst i poveznice
\definecolor{sivasvijetla}{HTML}{F2F4F5}
\definecolor{sivatamna}{HTML}{58595B}

\hypersetup{
  colorlinks=true,
  linkcolor=black,
  citecolor=fakplavatamna,
  urlcolor=fakplavatamna,
  pdftitle={\NaslovHR},
  pdfauthor={\Student}
}

% ---------------------------------------------------------------------
%  OSNOVNI PRIJELOM
% ---------------------------------------------------------------------
\onehalfspacing
\setlength{\parindent}{0pt}
\setlength{\parskip}{6pt}
\tolerance=2000
\emergencystretch=1.5em
\raggedbottom

\setlength{\headheight}{15pt}
\renewcommand{\headrule}{{\color{fakplava}\hrule height 0.6pt}}
\pagestyle{fancy}
\fancyhf{}
\fancyhead[L]{\small\sffamily\color{sivatamna}\NaslovKratki}
\fancyfoot[C]{\small\thepage}
\renewcommand{\headrulewidth}{0.6pt}
% ista glava i na stranicama na kojima počinje poglavlje
\fancypagestyle{plain}{%
  \fancyhf{}%
  \fancyhead[L]{\small\sffamily\color{sivatamna}\NaslovKratki}%
  \fancyfoot[C]{\small\thepage}%
  \renewcommand{\headrulewidth}{0.6pt}}

% ---------------------------------------------------------------------
%  NASLOVI POGLAVLJA I ODJELJAKA
% ---------------------------------------------------------------------
\titleformat{\chapter}[block]
  {\normalfont\sffamily\bfseries\fontsize{20}{24}\selectfont\raggedright}
  {\color{fakplava}\thechapter\hspace{0.6em}}
  {0pt}
  {}
  [\vspace{5pt}{\color{fakplava}\titlerule[2pt]}]
\titlespacing*{\chapter}{0pt}{0pt}{26pt}

\titleformat{\section}
  {\normalfont\sffamily\bfseries\fontsize{14}{17}\selectfont}
  {\color{fakplava}\thesection}{10pt}{}
\titlespacing*{\section}{0pt}{18pt}{6pt}

\titleformat{\subsection}
  {\normalfont\sffamily\bfseries\fontsize{12}{15}\selectfont}
  {\color{fakplava}\thesubsection}{8pt}{}
\titlespacing*{\subsection}{0pt}{14pt}{4pt}

\titleformat{\subsubsection}
  {\normalfont\sffamily\itshape\bfseries\fontsize{12}{15}\selectfont}
  {\thesubsubsection}{8pt}{}
\titlespacing*{\subsubsection}{0pt}{12pt}{4pt}

\setcounter{secnumdepth}{3}
\setcounter{tocdepth}{2}          % u Sadržaju se prikazuju razine 1.  i  1.1

% Numeracija slika, tablica i jednadžbi po poglavljima: 2.1, 2.2, ...
\counterwithin{figure}{chapter}
\counterwithin{table}{chapter}
\counterwithin{equation}{chapter}

% ---------------------------------------------------------------------
%  POTPISI (CAPTIONI)
% ---------------------------------------------------------------------
\captionsetup{
  font={small},
  labelfont={bf},
  labelsep=period,
  justification=centering,
  singlelinecheck=true,
  skip=8pt
}
\captionsetup[table]{position=top,skip=6pt}

% ---------------------------------------------------------------------
%  TABLICE
% ---------------------------------------------------------------------
\renewcommand{\arraystretch}{1.25}
\newcommand{\zaglavlje}[1]{\textbf{\textcolor{white}{#1}}}
\newcolumntype{L}[1]{>{\raggedright\arraybackslash}p{#1}}
\newcolumntype{C}[1]{>{\centering\arraybackslash}p{#1}}
\newcolumntype{R}[1]{>{\raggedleft\arraybackslash}p{#1}}

% ---------------------------------------------------------------------
%  KOD (LISTINZI)
% ---------------------------------------------------------------------
\lstdefinestyle{fak}{
  basicstyle=\ttfamily\footnotesize,
  backgroundcolor=\color{sivasvijetla},
  keywordstyle=\color{fakplavatamna}\bfseries,
  commentstyle=\color{sivatamna}\itshape,
  stringstyle=\color{sivatamna},
  numbers=left, numberstyle=\scriptsize\color{sivatamna}, numbersep=9pt,
  frame=leftline, framerule=2.5pt, rulecolor=\color{fakplava},
  framexleftmargin=26pt, xleftmargin=26pt,
  breaklines=true, showstringspaces=false,
  captionpos=b, aboveskip=10pt, belowskip=6pt,
  literate={č}{{\v{c}}}1 {ć}{{\'c}}1 {ž}{{\v{z}}}1 {š}{{\v{s}}}1 {đ}{{d}}1
}
\lstset{style=fak}
\renewcommand{\lstlistingname}{Programski isječak}

% Okvir s napomenom (koristi se u uputama, slobodno obriši)
\newcommand{\napomena}[1]{%
  \par\vspace{6pt}
  \noindent\begin{tabular}{@{}p{3pt}@{\hspace{10pt}}p{\dimexpr\textwidth-13pt\relax}@{}}
  \cellcolor{fakplava} & \small #1
  \end{tabular}\par\vspace{6pt}}

% Pomoćna naredba: ispiši drugi argument samo ako prvi nije prazan
\makeatletter
\newcommand{\akoNijePrazno}[2]{\begingroup\edef\@tempa{#1}%
  \ifx\@tempa\@empty\endgroup\else\endgroup #2\fi}
\makeatother

% Placeholder umjesto slike — zamijeni ga s \includegraphics
\newcommand{\mjestoZaSliku}[2][0.6\textwidth]{%
  \fcolorbox{fakplava}{sivasvijetla}{%
    \begin{minipage}[c][#2][c]{#1}
      \centering\sffamily\small\color{sivatamna}
      MJESTO ZA SLIKU\\[2pt]
      {\footnotesize zamijeni s \texttt{\textbackslash includegraphics}}
    \end{minipage}}}

% ---------------------------------------------------------------------
%  HRVATSKI NAZIVI
% ---------------------------------------------------------------------
\AtBeginDocument{%
  \renewcommand{\contentsname}{Sadržaj}
  \renewcommand{\listfigurename}{Popis slika}
  \renewcommand{\listtablename}{Popis tablica}
  \renewcommand{\bibname}{Literatura}
  \renewcommand{\figurename}{Slika}
  \renewcommand{\tablename}{Tablica}
  \renewcommand{\chaptername}{Poglavlje}
  \renewcommand{\appendixname}{Prilog}
}

% =====================================================================
%  BAZA LITERATURE — ugrađena u istu datoteku da sve ostane u jednom fileu.
%  Dodaj svoje unose ovdje. BibTeX ih sam sortira i formatira.
% =====================================================================
\begin{filecontents*}[overwrite]{\jobname.bib}
@book{knjiga2021,
  author    = {Prezime, Ime and Drugi, Autor},
  title     = {Naslov knjige: podnaslov ako postoji},
  publisher = {Naziv nakladnika},
  address   = {Zagreb},
  year      = {2021}
}
@article{clanak2023,
  author  = {Horvat, Ana and Kovačević, Marko},
  title   = {Naslov članka u znanstvenom časopisu},
  journal = {Naziv časopisa},
  volume  = {42},
  number  = {3},
  pages   = {115--129},
  year    = {2023},
  doi     = {10.1000/primjer.2023.42.115}
}
@inproceedings{konferencija2022,
  author    = {Novak, Ivan},
  title     = {Naslov rada s konferencije},
  booktitle = {Zbornik radova 15. međunarodne konferencije o nečemu},
  pages     = {44--51},
  address   = {Pula},
  year      = {2022}
}
@phdthesis{disertacija2020,
  author = {Babić, Petra},
  title  = {Naslov doktorske disertacije},
  school = {Sveučilište Jurja Dobrile u Puli, Fakultet informatike u Puli},
  year   = {2020}
}
@misc{mreznistvor2024,
  author       = {{Državni zavod za statistiku}},
  title        = {Naslov mrežne stranice ili dokumenta},
  howpublished = {\url{https://www.primjer.hr/dokument}},
  year         = {2024},
  note         = {pristupljeno 15. rujna 2026.}
}
@manual{norma2019,
  author       = {{International Organization for Standardization}},
  title        = {ISO 690:2010 Information and documentation — Guidelines for bibliographic references},
  organization = {ISO},
  address      = {Geneva},
  year         = {2010}
}
\end{filecontents*}

% =====================================================================
\begin{document}
% =====================================================================

% ---------------------------------------------------------------------
%  NASLOVNICA
% ---------------------------------------------------------------------
\begin{titlepage}
\thispagestyle{empty}
\centering
\sffamily
{\large \Sveuciliste\par}
{\large\bfseries \Fakultet\par}
{\normalsize \Odjel\par}

\vspace{10pt}
{\color{fakplava}\rule{\textwidth}{2pt}}

\vspace*{0.22\textheight}
{\large\bfseries\color{sivatamna} \VrstaRada\par}
\vspace{18pt}
{\fontsize{22}{27}\selectfont\bfseries \NaslovHR\par}
\vspace{10pt}
{\fontsize{13}{16}\selectfont\itshape\color{sivatamna} \NaslovEN\par}
\vfill

\begin{minipage}[t]{0.46\textwidth}
  \raggedright\normalsize
  \textbf{Student:}\\
  \Student\\
  {\small \JMBAG}
\end{minipage}\hfill
\begin{minipage}[t]{0.46\textwidth}
  \raggedleft\normalsize
  \textbf{Mentor:}\\
  \Mentor
  \akoNijePrazno{\Komentor}{\\[6pt]\textbf{Komentor:}\\ \Komentor}
\end{minipage}

\vspace{28pt}
{\color{fakplava}\rule{\textwidth}{2pt}}
\vspace{8pt}
{\normalsize \MjestoGodina\par}
\end{titlepage}

% ---------------------------------------------------------------------
%  IZJAVA
% ---------------------------------------------------------------------
\chapter*{Izjava o samostalnosti}
\addcontentsline{toc}{chapter}{Izjava o samostalnosti}

Izjavljujem da sam ovaj rad izradio/la samostalno, koristeći se znanjem stečenim tijekom studija i navedenom literaturom, te uz stručno vodstvo mentora.

\vspace{40pt}
\hfill\begin{minipage}{0.45\textwidth}
\centering
\rule{\linewidth}{0.4pt}\\
\small \Student
\end{minipage}

\napomena{Ako fakultet propisuje vlastiti obrazac izjave ili stranicu sa zadatkom rada, umetni ga ovdje umjesto ovog teksta.}

% ---------------------------------------------------------------------
%  ZAHVALA
% ---------------------------------------------------------------------
\chapter*{Zahvala}
\addcontentsline{toc}{chapter}{Zahvala}

Ovdje dolazi zahvala mentoru, obitelji, kolegama ili institucijama koje su podržale izradu rada. Zahvala nije obavezna — ako je ne želiš, obriši cijelo ovo poglavlje.

% ---------------------------------------------------------------------
%  SAŽETAK
% ---------------------------------------------------------------------
\chapter*{Sažetak}
\addcontentsline{toc}{chapter}{Sažetak}

Sažetak je samostalan tekst od 150 do 250 riječi koji se čita bez ostatka rada. U njemu u pravilu ide po jedna do dvije rečenice o svakome od sljedećeg: problem i motivacija, cilj rada, primijenjena metodologija, glavni rezultati i zaključak. U sažetku se ne citira literatura, ne koriste se kratice koje nisu odmah objašnjene i ne upućuje se na slike ili tablice iz rada.

\vspace{6pt}
\noindent\textbf{Ključne riječi:} prva ključna riječ; druga; treća; četvrta; peta

\vspace{26pt}
\noindent{\sffamily\bfseries\fontsize{20}{24}\selectfont Abstract}\par
\vspace{5pt}
\noindent{\color{fakplava}\rule{\textwidth}{2pt}}
\vspace{4pt}

The abstract is a direct translation of the Croatian summary — same length, same structure, same content. Have it proofread: this is the part of your thesis most likely to be read by someone outside your faculty.

\vspace{6pt}
\noindent\textbf{Keywords:} first keyword; second; third; fourth; fifth

% ---------------------------------------------------------------------
%  SADRŽAJ I POPISI
% ---------------------------------------------------------------------
\tableofcontents
\listoffigures
\listoftables

% =====================================================================
\chapter{Kako se koristi ovaj predložak}
% =====================================================================

Ovaj dokument je istovremeno predložak i uputa. Sve što vidiš — od naslovnice do popisa literature — već je oblikovano prema pravilima koja slijede. Tvoj posao je zamijeniti tekst uputa vlastitim sadržajem i ne dirati oblikovanje.

\napomena{Prije predaje obriši sva poglavlja s uputama (poglavlja 1 do 6) i zadrži samo strukturu. Popis za provjeru prije predaje nalazi se na kraju, u poglavlju~\ref{pog:predaja}.}

\section{Osnovna pravila prijeloma}

Sljedeće postavke su već ugrađene u predložak i ne mijenjaju se:

\begin{itemize}[leftmargin=*,itemsep=2pt,topsep=4pt]
  \item Format stranice A4, margine: lijevo 3~cm, desno 2{,}5~cm, gore i dolje 2{,}5~cm. Šira lijeva margina ostavlja mjesta za uvez.
  \item Osnovni tekst: Times New Roman 12~pt, prored 1{,}5, obostrano poravnanje.
  \item Naslovi: bezserifno pismo (Arial/Helvetica), broj poglavlja u boji fakulteta.
  \item Zaglavlje: skraćeni naslov rada uz tanku vodoravnu liniju, na svakoj stranici osim naslovnice.
  \item Broj stranice: sredina podnožja. Naslovnica se ne numerira.
  \item Odlomci se odvajaju razmakom, bez uvlake prvog retka.
\end{itemize}

\section{Jezik i stil}

Rad se piše u trećem licu ili bezličnim oblicima: \emph{,,mjerenja su provedena''}, a ne \emph{,,proveo sam mjerenja''}. Piše se u prošlom vremenu kad se opisuje što je napravljeno, a u prezentu kad se iznose činjenice i zaključci.

Kratice se pri prvom spominjanju pišu u punom obliku, a skraćenica se navodi u zagradi: \emph{umjetna inteligencija (UI)}. Nakon toga se koristi samo kratica. Decimalni se brojevi u hrvatskom pišu zarezom (3{,}14), a između broja i mjerne jedinice ide razmak (25~kg, 10~\%).

\section{Opseg rada}

Okvirni opseg je 30--40 stranica za završni i 50--70 stranica za diplomski rad, ne računajući priloge. Mjerodavan je pravilnik tvojeg fakulteta — ako se razlikuje od ovih brojeva, vrijedi pravilnik.

% =====================================================================
\chapter{Struktura rada i naslovi poglavlja}
% =====================================================================

\section{Redoslijed dijelova rada}

Rad se slaže ovim redoslijedom: naslovnica, izjava o samostalnosti, zahvala (nije obavezna), sažetak s ključnim riječima na hrvatskom i engleskom, sadržaj, popis slika, popis tablica, popis kratica (ako ih ima puno), glavni dio rada, literatura i prilozi.

Glavni dio rada ima ustaljenu strukturu: \textbf{Uvod} (kontekst, problem, cilj i struktura rada), \textbf{Pregled područja} (što se već zna i što su drugi napravili), \textbf{Metodologija} (kako si to napravio, dovoljno detaljno da netko može ponoviti), \textbf{Rezultati} (što je izašlo, bez tumačenja), \textbf{Rasprava} (što to znači i koja su ograničenja) i \textbf{Zaključak} (odgovor na cilj iz uvoda i prijedlog daljnjeg rada).

\section{Razine naslova}

Koriste se najviše tri razine numeriranih naslova:

\begin{itemize}[leftmargin=*,itemsep=2pt,topsep=4pt]
  \item \textbf{prva razina} — naslov poglavlja (ovdje: \emph{2 Struktura rada i naslovi poglavlja}); uvijek počinje na novoj stranici,
  \item \textbf{druga razina} — naslov odjeljka (ovdje: \emph{2.2 Razine naslova}); dijeli poglavlje na logičke cjeline,
  \item \textbf{treća razina} — naslov pododjeljka (ovdje: \emph{2.2.1}); koristi se samo kad odjeljak stvarno treba dodatno razlomiti.
\end{itemize}

Četvrta razina naslova znači da je strukturu potrebno pojednostaviti, a ne produbiti. Iza naslova uvijek slijedi tekst — nikad odmah novi naslov.

\subsection{Primjer pododjeljka}

Ovo je naslov treće razine. U Sadržaju se prikazuju prve tri razine, dublje se ne prikazuju.

\napomena{U LaTeX inačici naslove pišeš naredbama \texttt{\textbackslash chapter}, \texttt{\textbackslash section} i \texttt{\textbackslash subsection}. Numeraciju nikada ne upisuj ručno — LaTeX je generira sam, kao i Sadržaj.}

\section{Nabrajanja}

Nenumerirano nabrajanje koristi se kad redoslijed nije bitan:

\begin{itemize}[leftmargin=*,itemsep=2pt,topsep=4pt]
  \item prva stavka nabrajanja,
  \item druga stavka nabrajanja,
  \item posljednja stavka nabrajanja.
\end{itemize}

Numerirano nabrajanje koristi se kad redoslijed jest bitan, na primjer za korake postupka:

\begin{enumerate}[leftmargin=*,itemsep=2pt,topsep=4pt]
  \item prvi korak postupka,
  \item drugi korak postupka,
  \item treći korak postupka.
\end{enumerate}

Stavke počinju malim slovom i završavaju zarezom, a posljednja točkom — osim ako su stavke pune rečenice, u kojem slučaju svaka počinje velikim slovom i završava točkom.

% =====================================================================
\chapter{Slike, tablice i jednadžbe}
% =====================================================================

Svaka slika i svaka tablica mora biti numerirana, imati potpis i biti spomenuta u tekstu \emph{prije} nego se pojavi. Ako nešto nije spomenuto u tekstu, ne ide u rad.

\section{Slike}

Potpis slike ide \textbf{ispod} slike, centriran, veličine 10~pt. Numeracija je po poglavljima: prva slika u trećem poglavlju je Slika~3.1. Na sliku se u tekstu upućuje ovako: ,,Shema sustava prikazana je na slici~\ref{sl:primjer}.''

\begin{figure}[H]
  \centering
  \mjestoZaSliku[0.62\textwidth]{4.2cm}
  \caption{Kratak i opisan potpis slike koji se može razumjeti bez čitanja teksta}
  \label{sl:primjer}
\end{figure}

Ako slika nije tvoja, izvor se navodi na kraju potpisa, a izvor mora biti u popisu literature:

\begin{figure}[H]
  \centering
  \mjestoZaSliku[0.62\textwidth]{3.4cm}
  \caption{Potpis preuzete slike. Izvor: \citep{clanak2023}}
  \label{sl:preuzeta}
\end{figure}

Praktična pravila za slike: rezolucija najmanje 300~dpi za rasterske slike, a dijagrame i grafove izvozi kao vektorsku grafiku (PDF ili SVG) jer ostaju oštri pri svakom povećanju. Tekst na slici ne smije biti manji od 9~pt. Grafovi moraju imati označene osi s mjernim jedinicama. Ne koristi boju kao jedini nositelj informacije — rad se često ispisuje crno-bijelo, pa razlikuj krivulje i vrstom linije ili oznakom.

\napomena{U LaTeXu sliku umećeš umjesto sivog okvira iznad, naredbom:\newline
\texttt{\footnotesize\textbackslash includegraphics[width=0.7\textbackslash textwidth]\{slike/naziv.pdf\}}\newline
Naredba \texttt{\textbackslash label} mora doći \emph{poslije} \texttt{\textbackslash caption}, inače referenca pokazuje na krivi broj.}

\section{Tablice}

Potpis tablice ide \textbf{iznad} tablice. Numeracija je također po poglavljima. Tablice nemaju okomite crte — koriste se samo vodoravne linije, i to tri: iznad zaglavlja, ispod zaglavlja i na kraju tablice. Zaglavlje je u boji fakulteta s bijelim podebljanim tekstom.

\begin{table}[H]
  \centering
  \caption{Kratak potpis tablice koji objašnjava što se u njoj nalazi}
  \label{tab:primjer}
  \small
  \begin{tabular}{L{4.2cm} C{2.6cm} R{2.6cm} R{2.6cm}}
    \toprule
    \rowcolor{fakplava}
    \zaglavlje{Naziv uzorka} & \zaglavlje{Oznaka} & \zaglavlje{Masa (g)} & \zaglavlje{Udio (\%)} \\
    \midrule
    Prvi uzorak    & A-01 & 12{,}40 & 38{,}2 \\
    Drugi uzorak   & A-02 &  9{,}15 & 28{,}2 \\
    Treći uzorak   & B-01 & 10{,}90 & 33{,}6 \\
    \midrule
    \textbf{Ukupno} &     & \textbf{32{,}45} & \textbf{100{,}0} \\
    \bottomrule
  \end{tabular}
\end{table}

Brojevi se u stupcu poravnavaju udesno i imaju jednak broj decimala. Tekst se poravnava ulijevo. Mjerna jedinica ide u zaglavlje stupca, a ne uz svaki broj. Prazna ćelija označava se crticom (--), a ne ostavlja praznom, kako bi bilo jasno da podatak nedostaje namjerno.

Tablica koja se proteže preko više stranica mora ponoviti zaglavlje na svakoj stranici. Ako tablica ima više od desetak redaka i nije ključna za razumijevanje, premjesti je u prilog i u tekstu se pozovi na nju.

\section{Jednadžbe}

Jednadžbe se pišu u uređivaču jednadžbi, nikad kao slika i nikad kao obični tekst. Centrirane su, a broj jednadžbe ide desno u okruglim zagradama:

\begin{equation}
  \bar{x} = \frac{1}{n}\sum_{i=1}^{n} x_i
  \label{jed:aritm}
\end{equation}

Na jednadžbu se upućuje brojem: ,,aritmetička sredina računa se prema izrazu~(\ref{jed:aritm})''. Svi simboli moraju biti objašnjeni odmah nakon jednadžbe: $\bar{x}$ je aritmetička sredina, $n$ broj mjerenja, a $x_i$ vrijednost $i$-tog mjerenja. Varijable se pišu kurzivom, a mjerne jedinice i oznake funkcija uspravno.

\section{Programski kod}

Kraći isječci koda umeću se u tekst u monospace pismu s numeriranim recima. Dulji kod (više od tridesetak redaka) ide u prilog ili u javni repozitorij na koji se u tekstu uputi poveznicom.

\begin{lstlisting}[language=Python,caption={Primjer isječka koda s potpisom ispod},label={kod:primjer}]
def aritmeticka_sredina(vrijednosti):
    """Vraca aritmeticku sredinu liste brojeva."""
    if not vrijednosti:
        raise ValueError("Lista ne smije biti prazna.")
    return sum(vrijednosti) / len(vrijednosti)
\end{lstlisting}

% =====================================================================
\chapter{Citiranje i popis literature}
% =====================================================================

Svaka tvrdnja koja nije tvoja i nije općepoznata činjenica mora imati izvor. To vrijedi i za slike, tablice, podatke, definicije i programski kod preuzet s interneta.

\section{Citiranje u tekstu}

Predložak koristi \textbf{numerički sustav}: izvor se u tekstu označava brojem u uglatim zagradama, a broj odgovara mjestu izvora u popisu literature. Primjer jednog izvora \citep{clanak2023}, primjer više izvora odjednom \citep{knjiga2021,konferencija2022,disertacija2020}.

Broj se stavlja prije interpunkcije na kraju rečenice. Ako se navodi ime autora, ono ide u tekst, a broj odmah iza njega: ,,Horvat i Kovačević \citep{clanak2023} pokazali su da\ldots''

Doslovni navod dulji od jednog retka piše se pod navodnicima s obveznim brojem stranice, na primjer \citep[str.~118]{clanak2023}. Doslovne navode koristi štedljivo — u tehničkim radovima gotovo uvijek je bolje parafrazirati.

\napomena{Ako tvoj fakultet propisuje autor–godina sustav (APA, Harvard), u preambuli zamijeni \texttt{\textbackslash bibliographystyle\{unsrtnat\}} s \texttt{\textbackslash bibliographystyle\{apalike\}} i koristi \texttt{\textbackslash citep} i \texttt{\textbackslash citet}. Sve ostalo ostaje isto.}

\section{Popis literature}

Popis se nalazi na kraju rada, prije priloga, i sadrži \emph{isključivo} izvore koji su citirani u tekstu. U numeričkom sustavu izvori su poredani redoslijedom prvog pojavljivanja, a ne abecedno.

Oblik zapisa ovisi o vrsti izvora. U ugrađenoj bazi literature na početku ove datoteke nalaze se gotovi primjeri za knjigu, članak u časopisu, rad s konferencije, doktorsku disertaciju, mrežni izvor i normu — kopiraj onaj koji ti treba i zamijeni podatke.

Opći okvir za oblikovanje bibliografskih navoda daje norma ISO~690 \citep{norma2019}, a mrežni izvori navode se kao u primjeru \citep{mreznistvor2024}. Kod mrežnih izvora uvijek se navodi datum pristupa jer se sadržaj mijenja. Ako izvor ima DOI, navodi se DOI umjesto poveznice. Wikipedia, blogovi i forumi nisu prihvatljivi izvori za tvrdnje u radu — služe samo kao polazište za pronalazak pravih izvora.

\napomena{Redoslijed prevođenja u LaTeXu je: \texttt{pdflatex} $\rightarrow$ \texttt{bibtex} $\rightarrow$ \texttt{pdflatex} $\rightarrow$ \texttt{pdflatex}. Bez toga se popis literature ne pojavi, a na mjestu citata ostanu upitnici. Na Overleafu je dovoljno dva do tri puta pritisnuti \emph{Recompile}.}

\section{Alati i plagiranje}

Preporučuje se korištenje alata za upravljanje literaturom (Zotero, Mendeley, JabRef) jer sami generiraju BibTeX zapise i popis literature u traženom stilu.

Radovi se provjeravaju sustavom za otkrivanje plagijata. Parafraza znači da si tekst napisao vlastitim riječima i drugačijom strukturom rečenice, uz navođenje izvora — zamjena nekoliko riječi sinonimima i dalje je plagijat. Ako si u izradi rada koristio generativne alate umjetne inteligencije, to se navodi u skladu s pravilima fakulteta.

% =====================================================================
\chapter{Prilozi i popis kratica}
% =====================================================================

U priloge ide sve što je važno za provjeru rada, ali bi prekinulo čitanje glavnog teksta: dulji programski kod, sirovi mjerni podaci, ankete, sheme, tehnički crteži, potvrde.

Prilozi se označavaju slovima (Prilog A, Prilog B) i svaki počinje na novoj stranici s vlastitim naslovom. Slike i tablice unutar priloga numeriraju se s prefiksom priloga: Slika A.1, Tablica B.2. Svaki prilog mora biti spomenut u glavnom tekstu.

Ako u radu koristiš više od desetak kratica, dodaj popis kratica iza popisa tablica. Popis je abecedno poredan i ima oblik:

\begin{table}[H]
  \centering
  \small
  \begin{tabular}{L{2.6cm} L{10cm}}
    \toprule
    \rowcolor{fakplava}
    \zaglavlje{Kratica} & \zaglavlje{Puni naziv} \\
    \midrule
    API & Application Programming Interface \\
    DOI & Digital Object Identifier \\
    UI  & umjetna inteligencija \\
    \bottomrule
  \end{tabular}
\end{table}

% =====================================================================
\chapter{Prije predaje}
\label{pog:predaja}
% =====================================================================

Provjeri redom:

\begin{enumerate}[leftmargin=*,itemsep=3pt,topsep=4pt]
  \item Sva poglavlja s uputama (1--6) su obrisana.
  \item Na naslovnici su točni naslov, ime, JMBAG, mentor i datum.
  \item Sažetak i Abstract imaju 150--250 riječi i po pet ključnih riječi.
  \item Sadržaj, popis slika i popis tablica su osvježeni i brojevi stranica se poklapaju.
  \item Svaka slika i tablica ima potpis, broj i spomenuta je u tekstu.
  \item Svaka preuzeta slika, tablica i tvrdnja ima naveden izvor.
  \item Svaki izvor iz popisa literature citiran je u tekstu i obrnuto.
  \item Nema praznih naslova, ,,TODO'' bilješki ni neriješenih referenci (\texttt{??}).
  \item Pravopis i gramatika su provjereni, a Abstract pročitan od nekoga tko dobro zna engleski.
  \item Rad je izvezen u PDF i provjeren je izgled PDF-a, a ne samo izvorne datoteke.
  \item Datoteka je imenovana prema pravilima fakulteta, na primjer:\\
        \texttt{\small Prezime\_Ime\_diplomski.pdf}
\end{enumerate}

% =====================================================================
%  LITERATURA
% =====================================================================
\cleardoublepage
\addcontentsline{toc}{chapter}{Literatura}
\bibliographystyle{unsrtnat}
\bibliography{\jobname}

% ---------------------------------------------------------------------
%  ALTERNATIVA BEZ BibTeX-a
%  Ako ne želiš pokretati bibtex ili ti je potrebno strogo hrvatsko
%  oblikovanje navoda, zakomentiraj tri retka iznad i odkomentiraj
%  blok ispod. Redoslijed \bibitem-a određuje brojeve [1], [2], ...
% ---------------------------------------------------------------------
% \begin{thebibliography}{99}
% \bibitem{clanak2023} A. Horvat i M. Kovačević, ,,Naslov članka u znanstvenom časopisu'', \emph{Naziv časopisa}, sv. 42, br. 3, str. 115--129, 2023. DOI: 10.1000/primjer.2023.42.115.
% \bibitem{knjiga2021} I. Prezime i A. Drugi, \emph{Naslov knjige: podnaslov ako postoji}, 3. izd. Zagreb: Naziv nakladnika, 2021.
% \bibitem{konferencija2022} I. Novak, ,,Naslov rada s konferencije'', u \emph{Zbornik radova 15. međunarodne konferencije o nečemu}, Rijeka, 2022, str. 44--51.
% \bibitem{disertacija2020} P. Babić, ,,Naslov doktorske disertacije'', doktorska disertacija, Sveučilište u Rijeci, Naziv fakulteta, 2020.
% \bibitem{norma2019} ISO 690:2010 \emph{Information and documentation --- Guidelines for bibliographic references}, Geneva: ISO, 2010.
% \bibitem{mreznistvor2024} Državni zavod za statistiku, ,,Naslov mrežne stranice ili dokumenta'', 2024. [Mrežno]. Dostupno: \url{https://www.primjer.hr/dokument} [pristupljeno 15. rujna 2026.].
% \end{thebibliography}

% =====================================================================
%  PRILOZI
% =====================================================================
\appendix
\renewcommand{\thechapter}{\Alph{chapter}}

\chapter{Naziv prvog priloga}

Sadržaj priloga. Primjerice cjeloviti programski kod, tablica sirovih mjerenja ili obrazac ankete.

\chapter{Naziv drugog priloga}

Svaki prilog počinje na novoj stranici i ima svoj naslov.

\end{document}
